\begin{proofbox}
\textbf{方法一（乘法原理）}：
\begin{enumerate}
    \item 考虑集合 $A$ 中的每一个元素 $a_i$（$i=1,2,\dots,n$），在构造一个子集时，每个元素都有两种可能：被选入子集或不选入子集。
    \item 根据分步乘法计数原理，所有可能的选取方式共有 $\underbrace{2 \times 2 \times \cdots \times 2}_{n\text{个}} = 2^n$ 种。
    \item 每一种选取方式对应一个子集，且不同的选取方式对应不同的子集，故子集总数为 $2^n$。
\end{enumerate}
\end{proofbox}

\begin{proofbox}
\textbf{方法二（二项式定理）}：
\begin{enumerate}
    \item 按子集中所含元素的个数 $k$ 分类，$k$ 可取 $0,1,2,\dots,n$。
    \item 元素个数为 $k$ 的子集个数为组合数 $\mathrm{C}_n^k$。
    \item 因此子集总数为 $\sum_{k=0}^n \mathrm{C}_n^k = (1+1)^n = 2^n$（二项式定理）。
    \item 真子集即除去 $k=n$ 的情形，个数为 $2^n - \mathrm{C}_n^n = 2^n-1$；
    \item 非空子集即除去 $k=0$ 的情形，个数为 $2^n - \mathrm{C}_n^0 = 2^n-1$；
    \item 非空真子集即同时除去 $k=0$ 和 $k=n$，个数为 $2^n - \mathrm{C}_n^0 - \mathrm{C}_n^n = 2^n-2$。
\end{enumerate}
\end{proofbox}